\appendix
\section{Ficha Técnica de la Encuesta EAIMCS (INE Bolivia)}
\label{anexo:ficha_tecnica}

La Tabla \ref{tab:ficha_tecnica_eaimcs} resume las características metodológicas oficiales de la Encuesta a la Industria Manufacturera, Comercio y Servicios (EAIMCS 2017--2018), según los registros del catálogo oficial ANDA del Instituto Nacional de Estadística.

\begin{table}[htbp]
\centering
\caption{Ficha Técnica Metodológica de la Encuesta EAIMCS 2017--2018}
\label{tab:ficha_tecnica_eaimcs}
\footnotesize
\begin{tabularx}{\textwidth}{>{\bfseries\raggedright\arraybackslash}p{4.8cm} >{\raggedright\arraybackslash}X}
\toprule
\textbf{Parámetro Técnico} & \textbf{Especificación Oficial} \\
\midrule
Código Catálogo ANDA & \texttt{BOL-INE-EAIMCS-2017-2018} \\
Entidad Ejecutora & Instituto Nacional de Estadística (INE), Estado Plurinacional de Bolivia \\
Período de Referencia & Gestión Contable Anual 2017--2018 \\
Población Objetivo & Empresas formales pertenecientes a los estratos Mediana y Gran Empresa \\
Marco Muestral & Directorio Integrado de FUNDEMPRESA y Cuentas Nacionales ($N = 10,044$) \\
Tamaño Muestral Efectivo & 3,153 unidades productivas informantes \\
Diseño Muestral & No probabilístico, censal por cuotas en estratos medianos y grandes \\
Cobertura Geográfica & Nacional, 9 departamentos de Bolivia \\
Clasificador de Actividades & CAEB (Clasificación de Actividades Económicas de Bolivia, CIIU Rev. 4) \\
Marco Legal de Reserva & Decreto Ley N$^\circ$ 1405 (1976): Secreto y Confidencialidad Estadística \\
\bottomrule
\end{tabularx}
\vspace{2mm}
\raggedright\footnotesize
\textit{Nota.} Fuente: Instituto Nacional de Estadística (INE), Catálogo ANDA BOL-INE-EAIMCS-2017-2018. Los microdatos corresponden al operativo censal por cuotas en estratos de medianas y grandes empresas bajo resguardo de confidencialidad del Decreto Ley N$^\circ$ 1405.
\end{table}

\section{Matriz de Marco Lógico (MML) del Proyecto}
\label{anexo:marco_logico}

La Tabla \ref{tab:mml_resumen} presenta la síntesis formal de la Matriz de Marco Lógico que guió la ejecución del proyecto, detallando la lógica vertical, indicadores, medios de verificación y supuestos críticos.

\begin{table}[htbp]
\centering
\caption{Matriz Resumida de Marco Lógico del Proyecto}
\label{tab:mml_resumen}
\scriptsize
\setlength{\tabcolsep}{4pt}
\begin{tabularx}{\textwidth}{>{\bfseries\raggedright\arraybackslash}p{2.3cm} >{\raggedright\arraybackslash}X >{\raggedright\arraybackslash}p{3.5cm} >{\raggedright\arraybackslash}p{3.4cm}}
\toprule
\textbf{Nivel Jerárquico} & \textbf{Resumen Narrativo} & \textbf{Indicadores Verificables} & \textbf{Supuestos Críticos} \\
\midrule
\textbf{Fin (Impacto Esperado)} & Reducción de la brecha de recaudación fiscal e incremento en la confiabilidad de las mediciones macroeconómicas agregadas. & Disminución de la discrepancia residual promedio agregada y elevación de la consistencia del PIB sectorial. & Las entidades receptoras (SIN / INE) adoptan e institucionalizan formalmente el uso de la herramienta. \\
\addlinespace
\textbf{Propósito (Objetivo General)} & Identificación y contrastación oportuna de inconsistencias entre los ingresos declarados por medianas/grandes empresas y su rendimiento operativo estimado. & 1. $\text{MedAPE} \le 40\%$ y $R^2 \ge 0.60$ en 5-Fold Stratified CV.\newline 2. Cobertura empírica del IC al 90\% entre 85\% y 95\%. & Futuras rondas censales de EAIMCS o registros de facturación electrónica conservan compatibilidad estructural. \\
\addlinespace
\textbf{Componente 1 (Modelado Predictivo)} & Implementación de modelos analíticos e inteligentes de estimación de ingresos basados en capacidad productiva e insumos. & Pipeline campeón calibrado ($\hat{S} = 1.0401$, $R^2 = 0.7868$, $\text{MedAPE} = 36.20\%$) y matrices sectoriales normalizadas. & Microdatos de la EAIMCS representativos de los estratos medianos y grandes por actividad CAEB. \\
\addlinespace
\textbf{Componente 2 (Criterios de Selección)} & Automatización de la selección de contribuyentes para fiscalización basada en la brecha fiscal estimada. & 1. 100\% de unidades clasificadas en semáforos de riesgo (Bajo/Medio/Alto).\newline 2. Latencia del simulador $< 200$ ms. & Personal técnico cuenta con capacitación para interpretar el score como discrepancia y no como evasión tipificada. \\
\addlinespace
\textbf{Componente 3 (Parametrización Dinámica)} & Establecimiento de un esquema dinámico de parametrización y actualización de indicadores de control. & Alerta automática ante pruebas KS y Wasserstein ($D_{KS}, W_1$) con corrección de Bonferroni ($\alpha_{\text{ajustado}} = 0.004545$). & Se dispone de cargas periódicas de datos para que el monitoreo de deriva estadística tenga utilidad operativa. \\
\bottomrule
\end{tabularx}
\vspace{2mm}
\raggedright\footnotesize
\textit{Nota.} Estructura metodológica basada en los estándares del Enfoque del Marco Lógico (EML) de CEPAL / BID, articulando la lógica vertical con indicadores cuantitativos auditables.
\end{table}

\section{Diccionario de Variables Productivas Empleadas}
\label{anexo:diccionario_variables}

La Tabla \ref{tab:variables_modelo} especifica los predictores y la variable objetivo utilizados en el entrenamiento del modelo de regresión supervisada.

\begin{table}[htbp]
\centering
\caption{Diccionario de Predictores y Variable Objetivo en el Modelo Predictivo}
\label{tab:variables_modelo}
\footnotesize
\setlength{\tabcolsep}{4.5pt}
\begin{tabularx}{\textwidth}{l l >{\raggedright\arraybackslash}X l}
\toprule
\textbf{Variable} & \textbf{Sección INE} & \textbf{Descripción Técnica y Económica} & \textbf{Tipo de Dato} \\
\midrule
\texttt{target} & Sección 5 (\texttt{S00\_01\_A}) & Ingreso operativo anual total declarado (Bolivianos) & Numérico (Target) \\
\texttt{S01\_05\_A} & Sección 1 & Personal ocupado permanente y eventual total & Numérico \\
\texttt{S01\_03\_C} & Sección 1 & Remuneraciones y masa salarial total anual pagada (Bs) & Numérico \\
\texttt{S02\_09} & Sección 2 & Gasto anual total devengado en energía eléctrica (Bs) & Numérico \\
\texttt{S07\_09\_E} & Sección 7 & Valor histórico final de los activos fijos de capital (Bs) & Numérico \\
\texttt{S06\_06\_B} & Sección 6 & Valor de los inventarios finales de productos y materias (Bs) & Numérico \\
\texttt{S12\_01\_B} & Sección 12 & Capacidad física de almacenamiento construida ($m^2$) & Numérico \\
\texttt{S12\_02\_B} & Sección 12 & Capacidad física de almacenamiento total de predios ($m^2$) & Numérico \\
\texttt{n\_insumos} & Sección 10 & Cantidad de variedades distintas de materias primas & Numérico \\
\texttt{total\_valor\_co} & Sección 10 & Valor monetario total de compras de materias primas (Bs) & Numérico \\
\texttt{total\_valor\_uti}& Sección 10 & Valor monetario total consumido en producción fabril (Bs)& Numérico \\
\texttt{C2\_01} & Carátula & Departamento de localización geográfica (9 categorías) & Categórico \\
\texttt{macrosector} & Carátula & Sector de actividad económica normalizado (14 macrosectores)& Categórico \\
\bottomrule
\end{tabularx}
\vspace{2mm}
\raggedright\footnotesize
\textit{Nota.} Todas las variables monetarias y de escala física continuas son transformadas mediante $\ln(1+x)$ en el pipeline de ingeniería de características para corregir la asimetría distributiva. Las capacidades de almacenamiento ($S12\_01\_B$ y $S12\_02\_B$) forman parte del catálogo original de variables evaluadas, habiendo sido excluidas del entrenamiento final por alta heterogeneidad en las unidades de reporte según la Decisión D07.
\end{table}
